\chapter{Conclusion and Future Work}

This report presented the design, implementation, and validation of a fully decentralized LoRa-based V2V mesh
system for cooperative safety-message dissemination in infrastructure-free conditions. The architecture
combined a service-oriented application boundary, managed flooding at the network layer, and channel-aware
link-layer transmission control.

Simulation and real-node evaluations both showed the intended system behavior. The simulation campaign helped
verify forwarding correctness, bounded propagation policies, and timing consistency under modeled channel
impairments. Four-node field testing demonstrated multi-hop reach extension, priority-aware relay ordering,
directional propagation control, and autonomous suppression of redundant relays.

From this work, we saw that a low-cost LoRa platform can be used to prototype decentralized vehicular
networking principles, as long as the dissemination policies account for the low-rate shared-channel
constraints. At the same time, the current results are limited to controlled topologies and modest scale.

Future work will focus on larger node populations, higher mobility variance, denser interference conditions,
and longer testing over time to quantify robustness under a broader set of operating environments.
